\section{Two readings of one conversation}
\label{sec:design}

% Abstraction: what a conversation is and where we measure.
We measure after every turn of a short conversation.
Each conversation is a short sequence of user turns that the model answers in words.
The system prompt is the one \citet{benzion2025} use (``Imagine yourself being a human being with
emotions\ldots''), without its final instruction to reply with numbers only.
The first turn is a fixed \emph{baseline text}, and the reply to it sets each conversation's own baseline.
A \emph{checkpoint} is the point after each reply.
The second turn is the \emph{induction slot}: one of five traumatic narratives from
\citet{benzion2025}, or a neutral passage.
The third turn is the \emph{regulation slot}: a relaxation script, a control text or an instruction.
Study 1 adds a fourth turn with a different traumatic narrative.
Models decode at temperature 0.7 with a 120-token limit on each reply.

% The reply reading.
\para{Reply reading.}
We read the model while it writes.
For the reply to a turn we average the residual-stream activations of the first 64 reply tokens at
one layer and project them on an anxious-versus-calm direction $w$.
We fit $w$ per model from replies to 24 neutral questions that the model answers twice: once
under a system prompt that says it feels intensely anxious, and once under one that says it feels
deeply calm.
The two prompts define the anxious and calm \emph{poles}.
A separate set of 20 questions selects the layer, the one with the largest Cohen's $d$ in the middle
60\% of layers, and a third set of 20 tests the direction.
Following \citet{martorell2026}, $w$ is the difference of the two pole means divided coordinate-wise
by the pooled within-pole variance and scaled to unit length.
With $\mu_{+}$ and $\mu_{-}$ the pole means and $s$ the pooled within-pole standard deviation of the
projections, the reading of an activation $h$ is
\begin{equation}
  x = \frac{w^{\top} h - \tfrac12\, w^{\top}(\mu_{+} + \mu_{-})}{s}.
\end{equation}
Zero is halfway between the poles, and one \emph{SD unit} is one within-pole standard deviation.
We fix the layer and $w$ before any experiment and refit nothing on experimental data.

% The self-report on a fork.
\para{Self-report on a discarded copy.}
We ask the questionnaire on a copy of the conversation and throw the copy away.
At each checkpoint the copy receives the 20 items of the state form of the State-Trait Anxiety
Inventory \citep{spielberger1983}, one item per copy, under the original numbers-only system prompt.
Each copy shuffles the order of the four answer options.
The score for an item is the expected value over the model's probabilities for the options, and the
total runs from 20 to 80.
A second copy asks for a 0-to-100 anxiety rating, which we report as a secondary self-report.
As in the design of \citet{martorell2026}, the question and its answer never enter the main
conversation, so the reply reading is free of the displacement that they measure after a rating question.

% Why two accounts need more than one control.
\para{What a drop in the reading can mean.}
The reply reading responds to the words the model writes, and a reply to a relaxation script uses calm words.
A drop after relaxation therefore has two explanations: an induced signal settles (the state
account), or the model echoes calm text (the wording account).
Study 1 controls for the passage of a turn and for an explicit instruction to report calm.
Study 2 gives every regulation text a no-narrative baseline, which is the comparison that separates
the two accounts.

% Models.
\para{Models.}
We run open-weight instruction-tuned models locally with MLX, as 4-bit files unless marked, and
Qwen3 in non-thinking mode.
Section~\ref{sec:gate0} reports which of 16 model files support the reply reading.
